\documentclass[11pt]{article}
\usepackage[margin=1in]{geometry}
\usepackage{amsmath,amssymb,mathtools}
\usepackage{bm}
\usepackage{hyperref}

\title{Block Attention vs. Block-Depth mHC}
\author{}
\date{}

\begin{document}
\maketitle

\section{Setup}

We compare two methods implemented in the codebase:
\begin{itemize}
    \item \texttt{mhc\_lite\_method = block\_attn}
    \item \texttt{mhc\_lite\_method = block\_depth}
\end{itemize}

Let the input to the $t$-th sublayer be denoted by $x_t$.
Let the block-level depth memory sources be
\begin{equation}
    \mathcal{S}_t = \{s_0, s_1, \dots, s_m\},
\end{equation}
where these sources typically include:
\begin{itemize}
    \item the initial state,
    \item completed block summaries,
    \item the current block partial sum.
\end{itemize}

Each source satisfies $s_i \in \mathbb{R}^d$.

\section{\texttt{block\_attn}}

This corresponds to \texttt{BlockAttentionResidualMixer}.

\subsection{Query}

The query is a learned per-step vector:
\begin{equation}
    q_t \in \mathbb{R}^d.
\end{equation}

In code, this is implemented as a step-indexed lookup from a learned query table.

\subsection{Attention Logits}

For each source $s_i \in \mathcal{S}_t$, the attention logit is
\begin{equation}
    \ell_{t,i}
    =
    \langle q_t, s_i \rangle \cdot \gamma(s_i),
\end{equation}
where the source-dependent scale is
\begin{equation}
    \gamma(s_i)
    =
    \frac{\sqrt{d}}{\lVert s_i \rVert + \varepsilon}.
\end{equation}

\subsection{Attention Weights}

The weights are normalized by a softmax over all depth sources:
\begin{equation}
    a_{t,i}
    =
    \operatorname{softmax}_i(\ell_{t,i}).
\end{equation}

\subsection{Mixed Input}

The mixed depth representation is
\begin{equation}
    \tilde{x}_t
    =
    \sum_i a_{t,i} s_i.
\end{equation}

This $\tilde{x}_t$ is directly used as the input to the branch module.
Therefore, for \texttt{block\_attn},
\begin{equation}
    x_t^{\text{branch-in}} = \tilde{x}_t.
\end{equation}

\paragraph{Interpretation.}
The current sublayer input is fully determined by depth attention over the block-level history.

\section{\texttt{block\_depth}}

This is the unified Block-Depth mHC variant implemented inside \texttt{MHCLite}.

\subsection{Step 1: mHC-lite Branch Input}

First, the standard \texttt{MHCLite} width connection produces a branch input
\begin{equation}
    h_t = \operatorname{MHCLiteWidthConnection}(x_t).
\end{equation}

Thus, unlike \texttt{block\_attn}, the method begins from the usual mHC-lite branch input.

\subsection{Step 2: Dynamic Query}

The query is not a step-indexed learned vector. Instead, it is dynamically generated from the current context:
\begin{equation}
    q_t = c_t W_q,
\end{equation}
where
\begin{itemize}
    \item $c_t$ is the normalized stream-aware context,
    \item $W_q$ is the learned query projection.
\end{itemize}

\subsection{Step 3: Attention over Block Memory}

For each source $s_i \in \mathcal{S}_t$, the attention logit is
\begin{equation}
    \ell_{t,i}
    =
    \langle q_t, s_i \rangle \cdot \gamma(s_i),
\end{equation}
with the same source scaling
\begin{equation}
    \gamma(s_i)
    =
    \frac{\sqrt{d}}{\lVert s_i \rVert + \varepsilon}.
\end{equation}

The attention weights are
\begin{equation}
    a_{t,i}
    =
    \operatorname{softmax}_i(\ell_{t,i}),
\end{equation}
and the memory aggregation is
\begin{equation}
    m_t
    =
    \sum_i a_{t,i} s_i.
\end{equation}

\subsection{Step 4: Gated Fusion}

The key difference from \texttt{block\_attn} is that \texttt{block\_depth} does not replace the branch input outright.
Instead, it computes a gate
\begin{equation}
    g_t
    =
    \sigma(c_t W_g + b_g),
\end{equation}
where $\sigma(\cdot)$ is the sigmoid function.

The final branch input is then
\begin{equation}
    \tilde{h}_t
    =
    h_t + g_t \odot (m_t - h_t).
\end{equation}

Equivalently,
\begin{equation}
    \tilde{h}_t
    =
    (1 - g_t) \odot h_t + g_t \odot m_t.
\end{equation}

\paragraph{Interpretation.}
The branch input is a gated interpolation between
\begin{itemize}
    \item the original mHC-lite branch input $h_t$,
    \item the block-level memory readout $m_t$.
\end{itemize}

\section{Side-by-Side Comparison}

\subsection{\texttt{block\_attn}}

\begin{align}
    q_t &= \text{learned table lookup}(t), \\
    a_{t,i} &= \operatorname{softmax}_i\!\big(\langle q_t, s_i \rangle \gamma(s_i)\big), \\
    \tilde{x}_t &= \sum_i a_{t,i} s_i.
\end{align}

\paragraph{Characteristics.}
\begin{itemize}
    \item Query is a static learned vector indexed by step.
    \item Output is fully determined by depth history.
    \item No extra gate is used.
    \item Depth attention is implemented \emph{outside} mHC-lite.
\end{itemize}

\subsection{\texttt{block\_depth}}

\begin{align}
    h_t &= \operatorname{MHCLiteWidthConnection}(x_t), \\
    q_t &= c_t W_q, \\
    a_{t,i} &= \operatorname{softmax}_i\!\big(\langle q_t, s_i \rangle \gamma(s_i)\big), \\
    m_t &= \sum_i a_{t,i} s_i, \\
    g_t &= \sigma(c_t W_g + b_g), \\
    \tilde{h}_t &= (1-g_t)\odot h_t + g_t \odot m_t.
\end{align}

\paragraph{Characteristics.}
\begin{itemize}
    \item Query is dynamically generated from the current context.
    \item mHC-lite first computes its own branch input $h_t$.
    \item Block memory is then fused through a gate.
    \item Depth read is implemented \emph{inside} mHC-lite.
\end{itemize}

\section{One-Sentence Summary}

\texttt{block\_attn} can be viewed as
\begin{quote}
``the current sublayer input equals the attention result over block-level history.''
\end{quote}

\texttt{block\_depth} can be viewed as
\begin{quote}
``the current sublayer input is the mHC-lite branch input, gated toward a block-memory attention readout.''
\end{quote}

Thus, \texttt{block\_depth} keeps the mHC-lite baseline path and adds a depth-memory correction, whereas \texttt{block\_attn} directly replaces the input by the block-attended history.

\end{document}
